\documentclass{report}
\usepackage{amsmath,amssymb}
\setlength{\parindent}{0mm}
\setlength{\parskip}{1em}
\begin{document}
\begin{center}
\rule{6in}{1pt} \
{\large Dana Knoll \\
{\bf Computational Co-design of a moment-based scale-bridging algorithm }}

Los Alamos National Laboratory \\ PO BoX 1663 \\ MS B216 \\ Los Alamos \\ NM 87545
\\
{\tt nol@lanl.gov}\\
M. Daniels\\
C.N. Newman\\
H. Park\\
R.M. Rauenzahn\\
W. Taitano\\
J. Willert\end{center}

The foundations of the algorithmic scale-bridging ideas we will build on
reside in statistical mechanics, kinetic theory, and transport theory.
These ideas have matured in steady-state neutron transport for nuclear
reactors. In steady-state neutron transport, using this algorithm, the
transport (Boltzmann) equation would be referred to as the high order
(HO) problem and is six-dimensional (6-D). The low order (LO) problem is
constructed from combining the 0th and 1st phase-space moments of the
transport equation. This LO equation is only three-dimensional (3-D), and
thus is a reduced space representation of the HO problem. In the
appropriate limit, the LO problem asymptotes to the neutron diffusion
equation. This LO (moment-based) problem is used to accelerate the
convergence of a brute force transport source iteration via bridging the
transport and diffusion scales. This algorithm has characteristics of a
multigrid method. The LO problem lives in a reduced dimensional space
(3-D). It is an efficient way to relax the system-scale �diffusion�
features that naturally reside in the HO problem. The more expensive HO
problem (which lives in 6-D space) is only required to relax the
fine-scale transport features. The HO and LO problem are connected with
restriction and prolongation processes related to the moments.
Additionally, the LO problem is frequently on a coarser configuration
space mesh as compared to the HO problem. While this algorithmic idea is
frequently referred to as an accelerator for the transport (HO) solver,
it can also be viewed as a powerful physics-based scale-bridging
algorithm.

Development and use of similar moment-based acceleration methods can be
found in a variety of application areas such as neutron transport,
thermal radiation (photon) transport, kinetic plasma simulation, rarified
gas dynamics, and material science. We have ongoing efforts to advance
these algorithms in the areas of thermal radiation transport, kinetic
plasma simulation and neutron transport. In this talk we will draw the
majority of our example results from photon and neutron transport.

While much is known about these algorithms, and they are highly utilized
by some, further algorithm development and analysis remains. In this talk
we will provide an algorithmic overview with some application specific
details. Next we will discuss two important algorithmic issues. The first
issue is enforcing discrete consistency between the HO solver and the LO
solver. The second issue is related to challenges that arise when the HO
solver is stochastic, such as the use of Monte Carlo as the HO solver for
photon transport. At that point we have a hybrid algorithm with both
deterministic features and stochastic features. Finally, this algorithm
has a number of characteristics that will be beneficial in moving to
exascale computing. This general moment-based HO/LO algorithm has natural
heterogeneity and concurrency. We will discuss our ongoing computation
co-design efforts on emerging architectures to maximize the impact of
these algorithms for future multiphysics, multiscale simulation.


\end{document}
